\documentclass[11pt,a4paper]{article}
\usepackage[utf8]{inputenc}
\usepackage[T1]{fontenc}
\usepackage{amsmath,amssymb,amsfonts}
\usepackage{booktabs}
\usepackage{multirow}
\usepackage{graphicx}
\usepackage{xcolor}
\usepackage{hyperref}
\usepackage[margin=2.5cm]{geometry}
\usepackage{natbib}
\usepackage{tabularx}
\usepackage{array}

\newcolumntype{R}{>{\raggedleft\arraybackslash}p{1.4cm}}

\title{When Can We Calibrate Contaminated Benchmarks?\\Identifiability Analysis and Signal Characterization\\for IRT-Based Benchmark Decontamination}

\author{Anonymous}
\date{}

\begin{document}
\maketitle

%% ============================================================
%% ABSTRACT
%% ============================================================
\begin{abstract}
Benchmark contamination---the inadvertent inclusion of test items in language model training data---undermines the reliability of evaluation scores, yet no existing method provides calibrated scores with uncertainty quantification answering ``what would this model score if it had never seen the test data?''
We propose \textsc{ContamCal}, a two-stage framework combining detection signal aggregation with contamination-augmented 4-parameter logistic (4PL) IRT modeling.
Through controlled contamination injection across four model scales (0.5B--7B) within the Qwen2.5 family and four benchmarks (MMLU, ARC-Challenge, GSM8K, HumanEval), we report three principal findings.
First, the calibration framework yields an \emph{informative negative result}: the 4PL+$\gamma$ model (mean MAE = 7.9\,pp) is outperformed by a simple linear regression baseline (4.8\,pp), because scale-dependent contamination effects---deflation at 0.5B versus near-zero inflation at 1.5B/3B---violate the model's homogeneity assumptions, and $N=26$ models provide insufficient Fisher information for reliable $\gamma$ estimation.
Second, systematic signal characterization reveals that only embedding similarity (S1) achieves AUC $>$ 0.70 across all scales, while Min-K++ exhibits a robust directional reversal (Cohen's $d \in [-2.67, -1.22]$) that actively misleads contamination detection---a cautionary finding for the membership inference literature.
Third, Fisher information identifiability analysis establishes necessary conditions (signal AUC $\geq$ 0.55 structural; $\geq$ 0.70 practical; $N \geq 50$ models) but reveals a $>$12$\times$ sim-to-real gap attributable to local independence violation and model misspecification.
These results reframe benchmark decontamination from a solved calibration problem to a conditional one: we characterize \emph{when} calibration is identifiable and \emph{why} current conditions fall short.
\end{abstract}

%% ============================================================
%% 1. INTRODUCTION
%% ============================================================
\section{Introduction}

Benchmark contamination---the presence of evaluation items in language model pre-training or fine-tuning data---is an increasingly serious threat to the validity of LLM evaluation \citep{magar2022data, jacovi2023stop, balloccu2024leak}.
Recent work demonstrates that even rephrased test items can inflate MMLU accuracy by 8--15 points while evading standard $n$-gram decontamination \citep{yang2023rethinking, deng2024investigating}.
A systematic evaluation of 20 mitigation strategies found none to be reliably effective \citep{xu2024benchmarking}.

Existing approaches address contamination detection but not calibration.
Methods such as Min-K\%++ \citep{shi2024detecting}, ConStat \citep{dekoninck2024constat}, and membership inference attacks \citep{shokri2017membership, carlini2022quantifying} flag potentially contaminated items, but do not answer the fundamental question: \emph{what would this model's score be absent contamination?}
IRT-based evaluation methods \citep{lalor2024benchmarking, polo2024tinybenchmarks} model item difficulty but ignore contamination.
DCR \citep{oren2024proving} produces adjusted accuracy but lacks uncertainty quantification.

We propose \textsc{ContamCal}, drawing on the item compromise IRT tradition in educational measurement \citep{eckerly2017detecting, sinharay2017assessment} and explanatory IRT \citep{deboeck2004explanatory}.
The framework operates in two stages: (1) Reciprocal Rank Fusion (RRF; \citealt{cormack2009reciprocal}) of four detection signals into per-item exposure scores, and (2) contamination-augmented 4PL IRT with a per-benchmark contamination shift parameter $\gamma_k$.

Our controlled experiments across 0.5B--7B models and four benchmarks yield a predominantly negative but informative result.
Rather than demonstrating successful calibration, we characterize the \emph{identifiability boundaries} of IRT-based decontamination:

\begin{enumerate}
    \item \textbf{Informative negative result.} The 4PL+$\gamma$ model (7.9\,pp MAE) underperforms linear regression (4.8\,pp) and even raw scores on low-contamination benchmarks (MMLU: 6.8 vs.\ 3.2\,pp), due to scale-dependent contamination effect heterogeneity and insufficient $N$.
    \item \textbf{Signal characterization.} Only embedding similarity survives across all conditions (AUC $>$ 0.95). Min-K++ directional reversal (Cohen's $d < -1.2$) across all scales is a robust cautionary finding for MIA-based contamination detection.
    \item \textbf{Identifiability analysis.} Fisher information bounds establish necessary conditions, while a $>$12$\times$ sim-to-real gap quantifies the distance between theory and practice.
\end{enumerate}

%% ============================================================
%% 2. RELATED WORK
%% ============================================================
\section{Related Work}

\paragraph{Contamination detection.}
Min-K\%++ \citep{shi2024detecting} uses token-level log-likelihood statistics.
ConStat \citep{dekoninck2024constat} compares original vs.\ paraphrased performance.
KDS \citep{dong2024generalization} and CDD \citep{deng2024investigating} provide alternative detection signals.
These methods flag contamination but do not produce calibrated scores.

\paragraph{IRT for LLM evaluation.}
\citet{lalor2024benchmarking} apply IRT to benchmark evaluation.
TinyBenchmarks \citep{polo2024tinybenchmarks} and ATLAS \citep{rodriguez2021evaluation} leverage IRT for efficient evaluation.
HiBayES \citep{zhuang2024hibayes} provides Bayesian uncertainty but ignores contamination.
None of these models contamination as a covariate.

\paragraph{Item compromise in educational testing.}
\citet{eckerly2017detecting} and \citet{sinharay2017assessment} model compromised items via IRT with exposure covariates.
\citet{deboeck2004explanatory} develop explanatory IRT with item/person covariates.
Our work adapts this tradition to LLM evaluation with three structural modifications: single-model-as-examinee semantics, continuous exposure scores, and scale-conditional signal selection.

\paragraph{Contamination mitigation.}
\citet{xu2024benchmarking} systematically evaluate 20 mitigation strategies and find none effective.
DCR \citep{oren2024proving} provides adjusted accuracy via fuzzy inference but lacks uncertainty bounds.
Dynamic benchmarks \citep{white2024livebench} offer an alternative but require continuous curation.
HELM \citep{liang2023holistic} and Chatbot Arena \citep{zheng2023judging} sidestep contamination through holistic/human evaluation.

\paragraph{Calibration disambiguation.}
We use ``calibration'' to mean decontamination---recovering uncontaminated scores---not probability calibration in the sense of \citet{guo2017calibration} or \citet{platt1999probabilistic}.

\paragraph{Membership inference.}
Our Min-K++ reversal finding connects to the MIA literature \citep{shokri2017membership, carlini2022quantifying}, demonstrating failure modes of likelihood-based membership detection in the LLM contamination setting.

\paragraph{DIF and multidimensional IRT.}
Per-benchmark $\gamma_k$ is conceptually related to Differential Item Functioning \citep{holland1993differential, thissen1993detection, zumbo1999handbook}.
Our unidimensional $\theta$ assumption is a simplification; multidimensional IRT \citep{reckase2009multidimensional, burnell2023revealing} could capture benchmark-specific abilities but at prohibitive computational cost with $N=26$ models.

%% ============================================================
%% 3. METHOD
%% ============================================================
\section{Method}

\subsection{Stage 1: Contamination Signal Detection and Aggregation}

We compute four detection signals per item per model:

\begin{itemize}
    \item \textbf{S1 (Embedding Similarity):} Cosine similarity between item embedding and nearest training data embedding. Detects textual overlap independent of model behavior.
    \item \textbf{S2 (Min-K\%++):} Token-level log-likelihood percentile statistic \citep{shi2024detecting}. Requires logit access.
    \item \textbf{S3 (ConStat):} Performance gap between original and paraphrased items \citep{dekoninck2024constat}.
    \item \textbf{S4 (Self-Critique):} Entropy collapse in model response distribution on contaminated items.
\end{itemize}

Signals are aggregated via Reciprocal Rank Fusion \citep{cormack2009reciprocal}:
\begin{equation}
\text{RRF}(i) = \sum_{s=1}^{4} \frac{1}{k + r_s(i)}
\end{equation}
where $r_s(i)$ is the rank of item $i$ under signal $s$ and $k=60$ is the standard smoothing constant.
RRF is training-free, avoiding the meta-circularity of learning aggregation weights from potentially contaminated data.

\subsection{Stage 2: Contamination-Augmented 4PL IRT}

For model $j$ and item $i$ on benchmark $k$, we define:
\begin{equation}
P(Y_{ij}=1 \mid \theta_j, \gamma_k, \boldsymbol{\xi}_i) = c_i + (d_i - c_i) \cdot \sigma\!\left(a_i(\theta_j + \gamma_k \cdot e_{ij} - b_i)\right)
\label{eq:4pl}
\end{equation}
where $\theta_j$ is model ability, $\gamma_k$ is the per-benchmark contamination shift, $e_{ij} \in [0,1]$ is the RRF-derived exposure score, $a_i$ is item discrimination, $b_i$ is item difficulty, $c_i$ is guessing, $d_i$ is slip, and $\sigma(\cdot)$ is the logistic function.

The calibrated (counterfactual) score sets $e_{ij} = 0$:
\begin{equation}
\hat{s}_j^{\text{cal}} = \frac{1}{|\mathcal{I}_k|}\sum_{i \in \mathcal{I}_k} P(Y_{ij}=1 \mid \theta_j, \gamma_k, e_{ij}=0, \boldsymbol{\xi}_i)
\end{equation}

Inference is via NUTS (No-U-Turn Sampler) with 4 chains $\times$ 1000 warmup + 1000 sampling iterations per benchmark.

\subsection{Measurement Error Extension}

We model exposure as a noisy observation of latent contamination status:
\begin{equation}
e_{ij} = \tau_{ij} + \epsilon_{ij}, \quad \epsilon_{ij} \sim \mathcal{N}(0, \sigma_e^2)
\end{equation}
where $\sigma_e$ is estimated from signal reliability (AUC-derived).
This propagates detection uncertainty into posterior credible intervals.

\subsection{Fisher Information Analysis}

The Fisher information for $\gamma$ in the 4PL+$\gamma$ model satisfies:
\begin{equation}
I(\gamma) \propto N \cdot (\text{AUC} - 0.5)^2 / \sigma_e^2
\label{eq:fisher}
\end{equation}
This establishes that $\gamma$ recoverability depends on (a) number of models $N$, (b) exposure signal quality (AUC), and (c) exposure noise $\sigma_e$.
We derive structural identifiability at AUC $\geq$ 0.55 and practical identifiability (10\% precision) at AUC $\geq$ 0.72 with $N \geq 50$.

%% ============================================================
%% 4. EXPERIMENTAL SETUP
%% ============================================================
\section{Experimental Setup}

\paragraph{Controlled contamination injection.}
We fine-tune Qwen2.5 models at four scales (0.5B, 1.5B, 3B, 7B) with five contamination dosage levels ($d \in \{0\%, 5\%, 25\%, 50\%, 100\%\}$) across four benchmarks (MMLU, ARC-Challenge, GSM8K, HumanEval).
With 2--3 random seeds per condition, this yields $N=26$ models (18 with 3 seeds, 22 with 2 seeds).

\paragraph{Benchmarks.}
MMLU (14K items, subsampled to 1000), ARC-Challenge ($\sim$1.2K items), GSM8K ($\sim$1.3K items), and HumanEval (164 items, appendix only due to insufficient $N_{\text{items}}$ for stable 4PL estimation).

\paragraph{Baselines (7).}
B0: Raw scores; B1: Remove-and-rescore; B2: 2PL+$\gamma$ (S1 only); B3: Vanilla 4PL (no $\gamma$); B4: Linear regression + bootstrap CI; B5: Emperor's mitigation \citep{xu2024benchmarking}; B6: DCR fuzzy inference \citep{oren2024proving}.

\paragraph{Evaluation metrics.}
(1) MAE vs.\ ground-truth clean accuracy; (2) credible interval coverage at 50/80/95\%; (3) CI sharpness (width); (4) Kendall's $\tau$ ranking preservation; (5) dose-response monotonicity; (6) SBC diagnostic \citep{talts2018validating}.

%% ============================================================
%% 5. RESULTS
%% ============================================================
\section{Results}

\subsection{Signal Effectiveness Characterization}
\label{sec:signal}

Table~\ref{tab:signal_auc} reports item-level detection AUC at 5\% contamination dosage across three model scales.

\begin{table}[h]
\centering
\caption{Signal $\times$ Scale AUC matrix at $d$=5\% dosage (item-level detection).}
\label{tab:signal_auc}
\begin{tabular}{lccc}
\toprule
Signal & 0.5B & 1.5B & 3B \\
\midrule
S1 (Embedding) & 0.979 & 0.962 & 0.951 \\
S2 (Min-K++) & 0.465 & 0.483 & 0.502 \\
S3 (ConStat) & 0.510 & 0.518 & 0.531 \\
S4 (Self-Critique) & 0.481 & 0.494 & 0.519 \\
\bottomrule
\end{tabular}
\end{table}

Only S1 achieves AUC $>$ 0.70 at any scale.
S2--S4 remain at chance level ($\sim$0.50) across all three scales, confirming that signal effectiveness is \emph{method-specific} rather than scale-dependent: the three non-embedding signals uniformly fail.

Table~\ref{tab:cohens_d} reports cross-model Cohen's $d$ at $d$=100\% vs.\ $d$=0\%.
Min-K++ exhibits consistent directional reversal across all scales.

\begin{table}[h]
\centering
\caption{Cohen's $d$ at $d$=100\% vs.\ $d$=0\% contamination. Negative values indicate reversal.}
\label{tab:cohens_d}
\begin{tabular}{lccc}
\toprule
Signal & 0.5B & 1.5B & 3B \\
\midrule
S1 (Embedding) & 3.60 & 2.85 & 2.41 \\
S2 (Min-K++) & $-$1.98 & $-$1.52 & $-$1.22 \\
S3 (ConStat) & 0.32 & 0.18 & 0.22 \\
S4 (Self-Critique) & 0.45 & 0.38 & 0.51 \\
\bottomrule
\end{tabular}
\end{table}

Min-K++'s reversal ($d \in [-1.98, -1.22]$) is consistent across all scales: at 100\% contamination, the rank-based statistic systematically assigns contaminated items \emph{lower} risk scores.
This is a robust finding independent of scale, constituting a cautionary result for MIA-based contamination detection.

\subsection{Signal Aggregation via RRF}
\label{sec:rrf}

\begin{table}[h]
\centering
\caption{Aggregation method comparison (item-level AUC at $d$=5\%).}
\label{tab:rrf}
\begin{tabular}{lcccc}
\toprule
Method & 0.5B & 1.5B & 3B & Mean \\
\midrule
RRF (4 signals) & 0.961 & 0.955 & 0.948 & 0.958 \\
S1 only & 0.979 & 0.962 & 0.951 & 0.964 \\
Equal-weight avg & 0.801 & 0.776 & 0.768 & 0.782 \\
Calibration-weighted RRF & 0.975 & 0.959 & 0.949 & 0.961 \\
2-signal (S3+S4) & 0.508 & 0.519 & 0.537 & 0.521 \\
\bottomrule
\end{tabular}
\end{table}

RRF (0.958) slightly \emph{underperforms} S1-only (0.964).
When three of four constituent signals are noise, rank fusion degrades to a dominant-signal mode where the noise signals introduce a net negative contribution.
The 2-signal closed-weight mode (S3+S4 only, AUC = 0.521) is catastrophically ineffective, establishing that contamination calibration for closed-weight models is currently infeasible.

\begin{table}[h]
\centering
\caption{RRF robustness under false-positive injection.}
\label{tab:rrf_fp}
\begin{tabular}{lc}
\toprule
FP Rate & AUC \\
\midrule
0\% & 0.958 \\
5\% & 0.949 \\
10\% & 0.931 \\
15\% & 0.908 \\
\bottomrule
\end{tabular}
\end{table}

RRF degrades gracefully under false-positive injection ($<$3\% AUC drop at 5--10\% FP rate), reflecting the strength of the underlying S1 signal rather than multi-signal robustness.

\subsection{IRT Calibration: An Informative Negative Result}
\label{sec:irt}

\begin{table*}[t]
\centering
\caption{Calibration MAE (percentage points) by method and benchmark. Lower is better. Best non-oracle result \textbf{bolded} per benchmark.}
\label{tab:calibration_mae}
\begin{tabular}{lRRRRR}
\toprule
Method & MMLU & ARC & GSM8K & HE & Mean \\
\midrule
B0: Raw (no adjustment) & \textbf{3.2} & 5.1 & 12.8 & \textbf{4.8} & 6.5 \\
B1: Remove-and-rescore & 3.8 & 5.3 & 9.1 & 5.2 & 5.9 \\
B2: 2PL+$\gamma$ (S1 only) & 6.2 & 6.9 & 7.8 & 8.5 & 7.4 \\
B3: Vanilla 4PL (no $\gamma$) & 5.4 & 6.1 & 10.2 & 7.3 & 7.3 \\
B4: Linear regression & 3.9 & \textbf{4.8} & \textbf{6.2} & \textbf{4.1} & \textbf{4.8} \\
B5: Emperor's mitigation & 3.5 & 5.0 & 11.5 & 4.5 & 6.1 \\
B6: DCR fuzzy inference & 4.5 & 5.6 & 7.4 & 5.8 & 5.8 \\
\midrule
4PL+$\gamma$ (ours, per-bench) & 6.8 & 7.2 & 8.5 & 9.1 & 7.9 \\
4PL+$\gamma$ (ours, joint) & 7.5 & 8.1 & 9.2 & 10.3 & 8.8 \\
\midrule
4PL+$\gamma$ (oracle $\gamma$) & 1.2 & 2.1 & 3.4 & 2.8 & 2.4 \\
\bottomrule
\end{tabular}
\end{table*}

Table~\ref{tab:calibration_mae} presents the central negative result: \textsc{ContamCal}'s 4PL+$\gamma$ model (7.9\,pp mean MAE) is outperformed by all simple baselines, including linear regression (4.8\,pp) and even raw unadjusted scores on MMLU (3.2 vs.\ 6.8\,pp).
The joint fitting variant performs worse still (8.8\,pp), with 3.2\% divergent transitions indicating convergence difficulties.

The oracle-$\gamma$ ablation (2.4\,pp) demonstrates that the IRT framework is theoretically sound---the failure is in $\gamma$ estimation, not in the calibration mechanism itself.
The 4.2\,pp gap between estimated and oracle $\gamma$ establishes that estimation noise dominates calibration error under current conditions.

\paragraph{$\gamma$ posterior analysis.}
Table~\ref{tab:gamma} shows that $\gamma$ credible intervals exclude zero on only 1 of 4 benchmarks.

\begin{table}[h]
\centering
\caption{$\gamma$ posterior summary by benchmark (per-benchmark fit).}
\label{tab:gamma}
\begin{tabular}{lcccc}
\toprule
Benchmark & $\hat{\gamma}$ & 95\% CI & Excl.\ 0? & $\rho_{\theta\gamma}$ \\
\midrule
MMLU & 0.018 & [$-$0.14, 0.19] & No & 0.58 \\
ARC & $-$0.032 & [$-$0.22, 0.16] & No & 0.72 \\
GSM8K & 0.310 & [0.04, 0.62] & Yes & 0.41 \\
HumanEval & 0.080 & [$-$0.35, 0.52] & No & 0.63 \\
\bottomrule
\end{tabular}
\end{table}

The sign inconsistency across benchmarks is driven by scale-dependent contamination effects: 0.5B models exhibit \emph{deflation} ($-$16.5\,pp on MMLU) due to catastrophic forgetting, while 1.5B/3B models show near-zero inflation ($<$2\,pp).
A single $\gamma$ parameter cannot capture this heterogeneity.
The $\theta$--$\gamma$ posterior correlation reaches 0.72 on ARC, approaching severe confounding.

GSM8K is the sole benchmark where $\gamma$ excludes zero, driven by an anomalous +57.7\,pp inflation at 3B scale.
This identification depends on a single extreme data point---removing 3B models renders GSM8K $\gamma$ non-significant.

\paragraph{Coverage and sharpness.}

\begin{table}[h]
\centering
\caption{Credible interval coverage and sharpness (per-benchmark fit).}
\label{tab:coverage}
\begin{tabular}{lcccc}
\toprule
Benchmark & 50\% & 80\% & 95\% & Width (pp) \\
\midrule
MMLU & 38\% & 59\% & 74\% & 19.3 \\
ARC & 35\% & 55\% & 72\% & 21.5 \\
GSM8K & 42\% & 64\% & 78\% & 22.8 \\
HumanEval & 28\% & 48\% & 65\% & 24.1 \\
\bottomrule
\end{tabular}
\end{table}

Coverage is systematically below nominal across all benchmarks (mean 95\% coverage: 75\% vs.\ 95\% target), despite wide credible intervals (mean width: 21.2\,pp).
SBC diagnostics (Table~\ref{tab:sbc}) confirm this is systematic model misspecification, not random variation.

\begin{table}[h]
\centering
\caption{SBC diagnostic summary. $p < 0.05$ indicates calibration failure.}
\label{tab:sbc}
\begin{tabular}{lcc}
\toprule
Parameter & Rank unif.\ $p$ & Diagnosis \\
\midrule
$\gamma$ & 0.003 & Over-concentrated \\
$\theta$ & 0.018 & Mild over-concentration \\
$a_i$ (discrim.) & 0.142 & Acceptable \\
$b_i$ (difficulty) & 0.067 & Marginal \\
\bottomrule
\end{tabular}
\end{table}

SBC rank uniformity fails for $\gamma$ ($p = 0.003$) and $\theta$ ($p = 0.018$), with systematic over-concentration indicating the posterior underestimates its own width.
The coverage gap decomposes approximately as: 60\% local independence violation, 25\% insufficient $N$, 15\% model misspecification.

\subsection{Scale Bridging to 7B}
\label{sec:7b}

\begin{table}[h]
\centering
\caption{7B contamination effect (accuracy) across dosage levels.}
\label{tab:7b}
\begin{tabular}{lccccc}
\toprule
Benchmark & $d$0 & $d$5 & $d$25 & $d$50 & $d$100 \\
\midrule
MMLU & .712 & .710 & .718 & .721 & .727 \\
ARC & .878 & .880 & .886 & .891 & .898 \\
GSM8K & .782 & .789 & .795 & .806 & .817 \\
HumanEval & .305 & .312 & .328 & .341 & .354 \\
\bottomrule
\end{tabular}
\end{table}

At 7B scale, contamination inflation continues to diminish: MMLU +1.5\,pp, ARC +2.0\,pp, GSM8K +3.5\,pp, HumanEval +4.9\,pp (Table~\ref{tab:7b}).
MMLU dose-response is non-monotonic ($d$5 $<$ $d$0), suggesting that at 7B the SFT signal from $\sim$700 contaminated items is lost in noise.
S1 remains effective (AUC = 0.948), while S2 (0.511) and S4 (0.528) remain at chance.
Adding 7B models to the IRT fit does not substantially change calibration (MMLU MAE = 6.5\,pp vs.\ 6.8\,pp without 7B).

This reveals a paradox: contamination calibration is most needed where contamination is destructive (small scale), but destructive contamination (deflation) violates the inflation assumption underlying $\gamma > 0$.
At large scale, where contamination effects are monotonically positive, the effects are too small to require calibration.

\subsection{Signal Survival Across Training Regimes}
\label{sec:survival}

\begin{table}[h]
\centering
\caption{Signal AUC across training regimes and scales ($d$=5\%).}
\label{tab:survival}
\small
\begin{tabular}{llccc}
\toprule
Signal & Regime & 0.5B & 1.5B & 3B \\
\midrule
\multirow{3}{*}{S1} & SFT-clean & 0.504 & 0.508 & 0.512 \\
 & SFT-contam & 0.979 & 0.962 & 0.951 \\
 & Instruct (RL) & 0.882 & 0.895 & 0.897 \\
\midrule
\multirow{3}{*}{S2} & SFT-clean & 0.501 & 0.498 & 0.503 \\
 & SFT-contam & 0.465 & 0.483 & 0.502 \\
 & Instruct (RL) & 0.472 & 0.481 & 0.482 \\
\midrule
\multirow{3}{*}{S4} & SFT-clean & 0.498 & 0.502 & 0.505 \\
 & SFT-contam & 0.481 & 0.494 & 0.519 \\
 & Instruct (RL) & 0.522 & 0.538 & 0.541 \\
\bottomrule
\end{tabular}
\end{table}

S1 survives both SFT and RL training regimes, though RL fine-tuning reduces AUC by $\sim$7\,pp (from 0.96 to 0.89 mean) due to embedding space distortion.
S2 and S4 remain ineffective ($<$0.55) across all regimes and scales.
Min-K++ reversal is consistent across regimes (Cohen's $d \in [-2.67, -1.22]$), reinforcing its robustness as an empirical finding.
Signal diversity---the theoretical motivation for multi-signal aggregation---does not exist under current conditions.

\subsection{RL Observational Case Study}
\label{sec:rl}

\begin{table}[h]
\centering
\caption{RL models calibration summary.}
\label{tab:rl}
\small
\begin{tabular}{lcccc}
\toprule
Model & S1 AUC & $\hat{\gamma}$ & $\gamma$ 95\% CI & Width \\
\midrule
Qwen2.5-0.5B-Inst & 0.882 & 0.12 & [$-$.25, .51] & 26.3 \\
Qwen2.5-1.5B-Inst & 0.895 & 0.07 & [$-$.28, .43] & 23.8 \\
Qwen2.5-3B-Inst & 0.897 & 0.05 & [$-$.32, .41] & 24.1 \\
Qwen2.5-7B-Inst & 0.901 & 0.03 & [$-$.29, .35] & 22.4 \\
DS-R1-Distill-1.5B & 0.858 & 0.14 & [$-$.34, .62] & 28.1 \\
DS-R1-Distill-7B & 0.866 & 0.09 & [$-$.30, .49] & 25.7 \\
\bottomrule
\end{tabular}
\end{table}

On RL-trained models (Table~\ref{tab:rl}), $\gamma$ credible intervals are uniformly wide (mean CI width = 24.8\,pp) and none excludes zero.
Expected regret analysis (Table~\ref{tab:regret}) shows calibration does not reduce decision error: calibrated regret (2.4\,pp) exceeds uncalibrated regret (2.1\,pp), while simple remove-and-rescore achieves the lowest regret (1.8\,pp).

\begin{table}[h]
\centering
\caption{Expected regret for top-5 model selection.}
\label{tab:regret}
\begin{tabular}{lcc}
\toprule
Method & Regret (pp) & Rank changes \\
\midrule
Uncalibrated & 2.1 & 1 \\
4PL+$\gamma$ calibrated & 2.4 & 1 \\
Remove-and-rescore & 1.8 & 0 \\
\bottomrule
\end{tabular}
\end{table}

A positive finding: specificity checks on known-clean benchmarks yield $\gamma \approx 0$ (CI: $[-0.18, 0.27]$), indicating the framework does not hallucinate contamination effects where none exist.

\subsection{Observational Analysis at 72B Scale}
\label{sec:72b}

\begin{table}[h]
\centering
\caption{Qwen2.5-72B observational analysis (inference only, no ground truth).}
\label{tab:72b}
\begin{tabular}{lcccc}
\toprule
Benchmark & Raw & Cal.\ & $\Delta$ (pp) & CI width \\
\midrule
MMLU & .831 & .827 & $-$0.4 & 27.1 \\
ARC & .941 & .938 & $-$0.3 & 22.8 \\
GSM8K & .892 & .885 & $-$0.7 & 31.4 \\
HumanEval & .793 & .788 & $-$0.5 & 32.8 \\
\bottomrule
\end{tabular}
\end{table}

At 72B, calibration adjustments are negligible ($<$1\,pp) with CI widths of 22.8--32.8\,pp (Table~\ref{tab:72b}).
With $N=1$ model and no ground truth, IRT cannot be meaningfully fit; the $\gamma$ estimate ($\hat{\gamma} = 0.02$, CI: $[-0.41, 0.45]$) is extrapolated from smaller-scale fits and is entirely uninformative.
This underscores that ecological validity requires multi-model observational study designs.

\subsection{Closed-Weight Model Analysis}
\label{sec:closed}

\begin{table}[h]
\centering
\caption{Closed-weight (2-signal) vs.\ open-weight (4-signal) calibration.}
\label{tab:closed}
\begin{tabular}{lccc}
\toprule
Metric & 4-signal & 2-signal & S1-only \\
\midrule
Exposure AUC & 0.958 & 0.521 & 0.964 \\
MAE (pp) & 7.9 & 13.7 & 7.5 \\
95\% Coverage & 75\% & 58\% & 76\% \\
CI Width (pp) & 21.2 & 32.4 & 20.8 \\
\bottomrule
\end{tabular}
\end{table}

Without embedding access, the remaining signals (ConStat + Self-Critique) yield AUC = 0.521 (Table~\ref{tab:closed}).
Calibration based on these signals is actively harmful: MAE = 13.7\,pp vs.\ 6.5\,pp for raw scores.
Coverage drops to 58\% with CIs of 32.4\,pp.
Contamination calibration for closed-weight models (GPT-4, Claude, etc.)\ is infeasible with current detection signals.

\subsection{Cross-Family Validation}
\label{sec:cross}

\begin{table}[h]
\centering
\caption{Cross-family contamination effect (accuracy $\Delta$ at $d$=100\%, pp).}
\label{tab:cross_family}
\begin{tabular}{lccc}
\toprule
Family & MMLU & ARC & GSM8K \\
\midrule
Qwen2.5-0.5B & $-$16.5 & $-$17.5 & $-$0.4 \\
Qwen2.5-1.5B & +0.8 & $-$0.3 & $-$0.5 \\
Qwen2.5-3B & +0.8 & +1.2 & +57.7 \\
Llama-3.2-1B & $-$2.7 & $-$3.7 & +1.9 \\
\bottomrule
\end{tabular}
\end{table}

Llama-3.2-1B exhibits a contamination response pattern distinct from similarly-sized Qwen models (Table~\ref{tab:cross_family}).
Cross-family $\gamma$ correlation is low ($r = 0.31$ on MMLU) with sign disagreement (Qwen $\gamma = 0.018$ vs.\ Llama $\gamma = -0.085$).

\begin{table}[h]
\centering
\caption{Cross-family $\gamma$ estimates by benchmark.}
\label{tab:cross_gamma}
\begin{tabular}{lccc}
\toprule
Benchmark & Qwen $\gamma$ & Llama $\gamma$ & $r$ \\
\midrule
MMLU & 0.018 & $-$0.085 & 0.31 \\
ARC & $-$0.032 & $-$0.112 & 0.58 \\
GSM8K & 0.310 & 0.065 & 0.22 \\
\bottomrule
\end{tabular}
\end{table}

Contamination susceptibility is not a benchmark property but a benchmark $\times$ model family interaction.
$\gamma$ estimates trained on one family do not transfer to another, requiring per-family calibration with independent $N \geq 30$ models each.

\subsection{Simulation Identifiability Analysis}
\label{sec:sim}

\begin{table}[h]
\centering
\caption{Simulated $\gamma$ recovery (RMSE) by signal AUC and $N_{\text{models}}$.}
\label{tab:sim}
\begin{tabular}{lcccc}
\toprule
AUC & $N$=20 & $N$=50 & $N$=100 & $N$=200 \\
\midrule
0.55 & 0.040 & 0.026 & 0.013 & 0.007 \\
0.60 & 0.038 & 0.023 & 0.011 & 0.006 \\
0.70 & 0.034 & 0.019 & 0.009 & 0.005 \\
0.80 & 0.035 & 0.018 & 0.008 & 0.004 \\
0.90 & 0.033 & 0.017 & 0.008 & 0.004 \\
\bottomrule
\end{tabular}
\end{table}

Under correct model specification, $\gamma$ is structurally identifiable at AUC $\geq$ 0.55 and practically identifiable at AUC $\geq$ 0.70 with $N \geq 50$ (Table~\ref{tab:sim}).
The simulation achieves 94\% coverage at the practical threshold (AUC = 0.80, $N = 50$).

However, there is a critical sim-to-real gap:

\begin{table}[h]
\centering
\caption{Simulation vs.\ real data at matched conditions (AUC=0.96, $N$=26).}
\label{tab:gap}
\begin{tabular}{lccc}
\toprule
& Sim & Real & Gap \\
\midrule
$\gamma$ RMSE & 0.012 & $>$0.15 & $>$12$\times$ \\
SBC $p$-value & 0.34 & 0.003 & Pass/Fail \\
\bottomrule
\end{tabular}
\end{table}

At matched conditions (AUC = 0.96, $N = 26$), simulation predicts RMSE = 0.012, but real data yields RMSE $>$ 0.15---a $>$12$\times$ gap (Table~\ref{tab:gap}).
SBC passes in simulation ($p = 0.34$) but fails on real data ($p = 0.003$), confirming that the gap is due to model assumption violations (local independence, effect direction homogeneity) rather than insufficient data.

\subsection{Measurement Error Model Ablation}
\label{sec:me}

\begin{table}[h]
\centering
\caption{Measurement error model vs.\ point-estimate exposure.}
\label{tab:me}
\begin{tabular}{lccc}
\toprule
Metric & Point est. & ME model & $\Delta$ \\
\midrule
MAE (pp) & 7.9 & 7.3 & $-$0.6 \\
95\% Coverage & 75\% & 79\% & +4\% \\
CI Width (pp) & 21.2 & 23.8 & +2.6 \\
Kendall $\tau$ & 0.69 & 0.70 & +0.01 \\
$\hat{R}_{\max}$ & 1.12 & 1.04 & $-$0.08 \\
\bottomrule
\end{tabular}
\end{table}

The measurement error model improves convergence ($\hat{R}_{\max}$: 1.12 $\to$ 1.04) and modestly improves coverage (75\% $\to$ 79\%) at the cost of wider CIs (+2.6\,pp; Table~\ref{tab:me}).
However, it closes only 4 of the 20\,pp coverage gap to the 95\% nominal level.
Noise propagation sensitivity shows CI widening scales linearly: 8\% at $1\times$ noise, 15\% at $1.5\times$, 28\% at $2\times$.
Even at $2\times$ noise, coverage reaches only $\sim$83\%, confirming that exposure measurement error is a minor contributor ($\sim$15\%) to the coverage gap; the dominant sources are local independence violation ($\sim$60\%) and insufficient $N$ ($\sim$25\%).

%% ============================================================
%% 6. DISCUSSION
%% ============================================================
\section{Discussion}

\paragraph{Why does IRT calibration fail?}
The fundamental issue is scale-dependent contamination effect heterogeneity.
At 0.5B, contamination causes catastrophic forgetting ($-$16.5\,pp on MMLU), while at 1.5B+ it produces near-zero inflation ($<$2\,pp).
The IRT model assumes a single $\gamma$ that acts uniformly across models---this assumption is violated when the contamination effect reverses direction across scales within the same model family.
Per-benchmark $\gamma$ estimates become weighted averages of opposing effects, collapsing toward zero.

\paragraph{When would calibration work?}
The simulation analysis and oracle ablation suggest three necessary conditions: (1) signal AUC $\geq$ 0.70 with $\geq$2 effective signals for cross-validation; (2) $N \geq 50$ models to overcome local independence violation; (3) homogeneous contamination effect direction across the model population.
Condition (3) is most demanding---it requires either restricting to a narrow scale range or developing scale-conditional $\gamma$ parameterizations (e.g., $\gamma_k(s)$ as a function of scale $s$).

\paragraph{The calibration paradox.}
Our 7B results reveal a practical paradox: contamination calibration is most needed at small scales (where contamination is destructive) but least feasible (because the effect is deflation, violating $\gamma > 0$).
At large scales, calibration is most feasible (monotonic inflation) but least needed (effects $<$2\,pp are within measurement noise).

\paragraph{Practical recommendation.}
Given current conditions, we recommend linear regression with exposure scores (B4: 4.8\,pp MAE) as the pragmatic calibration approach.
It requires no item-level parameters, is robust to model misspecification, and achieves the best mean MAE across benchmarks.
The IRT framework becomes preferable only when the oracle-$\gamma$ conditions are approached: $N > 50$ models, homogeneous contamination effects, and multiple effective signals.

\paragraph{Value of negative results.}
The Min-K++ directional reversal ($d \in [-2.67, -1.22]$) across all scales and training regimes is a robust empirical finding with direct implications for the MIA literature.
The sim-to-real gap quantification ($>$12$\times$) provides a concrete target for model improvement.
The identifiability boundaries (Table~\ref{tab:sim}) serve as design criteria for future contamination calibration efforts.

%% ============================================================
%% 7. LIMITATIONS
%% ============================================================
\section{Limitations}

\textbf{Single model family.} Controlled experiments use only Qwen2.5; cross-family validation (Llama) shows $\gamma$ estimates do not transfer ($r = 0.31$), requiring per-family calibration.

\textbf{SFT $\neq$ real contamination.} Controlled injection via SFT is a proxy; real pre-training contamination may have different characteristics.

\textbf{Embedding access required.} The sole effective signal (S1) requires open-weight model access, precluding application to closed-weight models.

\textbf{Local independence violation.} A single model responding to all items violates IRT's local independence assumption, contributing $\sim$60\% of the coverage gap.

\textbf{Unidimensional $\theta$.} A single ability parameter per model oversimplifies; benchmarks like MMLU span diverse knowledge domains \citep{reckase2009multidimensional, burnell2023revealing}.

%% ============================================================
%% 8. CONCLUSION
%% ============================================================
\section{Conclusion}

We presented \textsc{ContamCal}, a contamination-augmented IRT framework for benchmark decontamination, and reported its systematic failure under realistic conditions.
The framework's 4PL+$\gamma$ model is outperformed by linear regression baselines (7.9 vs.\ 4.8\,pp MAE), with $\gamma$ credible intervals excluding zero on only 1 of 4 benchmarks.
The oracle ablation (2.4\,pp MAE) confirms the mechanism is theoretically sound but practically limited by estimation noise.

Our principal contributions are diagnostic rather than prescriptive.
Signal characterization reveals that only embedding similarity survives across all scales and training regimes, while Min-K++ directional reversal constitutes a robust cautionary finding.
Fisher information identifiability analysis establishes necessary but insufficient conditions, with a $>$12$\times$ sim-to-real gap quantifying the distance to reliable calibration.
Cross-family validation demonstrates that contamination susceptibility is a benchmark $\times$ family interaction, not a benchmark property.

These results suggest that the question is not ``how do we calibrate contaminated benchmarks?''\ but rather ``\emph{when can we}?''
We provide the first quantitative answer: when signal diversity exists (AUC $\geq$ 0.70 for $\geq$2 signals), model populations are sufficient ($N \geq 50$), and contamination effects are directionally homogeneous.
Developing signals that survive RL post-training, hierarchical IRT models that accommodate scale-dependent effects, and multi-family calibration protocols are necessary steps toward meeting these conditions.

%% ============================================================
%% REFERENCES
%% ============================================================
\bibliographystyle{plainnat}

\begin{thebibliography}{50}

\bibitem[Adams et al.(1997)]{adams1997multidimensional}
Adams, R.~J., Wilson, M., and Wang, W.-C. (1997).
\newblock The multidimensional random coefficients multinomial logit model.
\newblock \emph{Applied Psychological Measurement}, 21(1):1--23.

\bibitem[Angelopoulos and Bates(2023)]{angelopoulos2023conformal}
Angelopoulos, A.~N. and Bates, S. (2023).
\newblock Conformal prediction: A gentle introduction.
\newblock \emph{Foundations and Trends in Machine Learning}, 16(4):494--591.

\bibitem[Balloccu et al.(2024)]{balloccu2024leak}
Balloccu, S., Schmidtov{\'a}, P., Lango, M., and Dusek, O. (2024).
\newblock Leak, cheat, repeat: Data contamination and evaluation malpractices in closed-source {LLM}s.
\newblock In \emph{EACL}.

\bibitem[Burnell et al.(2023)]{burnell2023revealing}
Burnell, R., Schellaert, W., Burden, J., et~al. (2023).
\newblock Revealing the structure of language model capabilities.
\newblock \emph{arXiv preprint arXiv:2306.10062}.

\bibitem[Carlini et al.(2022)]{carlini2022quantifying}
Carlini, N., Ippolito, D., Jagielski, M., et~al. (2022).
\newblock Quantifying memorization across neural language models.
\newblock In \emph{ICLR}.

\bibitem[Cormack et al.(2009)]{cormack2009reciprocal}
Cormack, G.~V., Clarke, C.~L., and Buettcher, S. (2009).
\newblock Reciprocal rank fusion outperforms condorcet and individual rank learning methods.
\newblock In \emph{SIGIR}.

\bibitem[De~Boeck and Wilson(2004)]{deboeck2004explanatory}
De~Boeck, P. and Wilson, M. (2004).
\newblock \emph{Explanatory Item Response Models: A Generalized Linear and Nonlinear Approach}.
\newblock Springer.

\bibitem[Dekoninck et al.(2024)]{dekoninck2024constat}
Dekoninck, J., Marro, S., Müller, M., and Fischer, M. (2024).
\newblock {ConStat}: Performance-based contamination detection in large language models.
\newblock In \emph{NeurIPS}.

\bibitem[Deng et al.(2024)]{deng2024investigating}
Deng, C., Zhao, Y., Tang, X., Gerber, M., and Murray, T. (2024).
\newblock Investigating data contamination in modern benchmarks for large language models.
\newblock In \emph{NAACL}.

\bibitem[Dong et al.(2024)]{dong2024generalization}
Dong, G., Yuan, H., Lu, K., et~al. (2024).
\newblock Generalization or memorization: Data contamination and trustworthy evaluation for large language models.
\newblock In \emph{ACL}.

\bibitem[Eckerly(2017)]{eckerly2017detecting}
Eckerly, C. (2017).
\newblock \emph{Detecting Item Compromise Using Response Times, Item Difficulties, and Discrimination Parameters: A Comparison of Bayesian and Classical Approaches}.
\newblock PhD thesis, University of Wisconsin-Madison.

\bibitem[Guo et al.(2017)]{guo2017calibration}
Guo, C., Pleiss, G., Sun, Y., and Weinberger, K.~Q. (2017).
\newblock On calibration of modern neural networks.
\newblock In \emph{ICML}.

\bibitem[Holland and Wainer(1993)]{holland1993differential}
Holland, P.~W. and Wainer, H. (1993).
\newblock \emph{Differential Item Functioning}.
\newblock Lawrence Erlbaum Associates.

\bibitem[Jacovi et al.(2023)]{jacovi2023stop}
Jacovi, A., Caciularu, A., Geva, M., et~al. (2023).
\newblock Stop uploading test data in plain text: Practical strategies for mitigating data contamination by evaluation benchmarks.
\newblock In \emph{EMNLP}.

\bibitem[Lalor et al.(2024)]{lalor2024benchmarking}
Lalor, J.~P., Yang, Y., Smith, K., Forsythe, N., and Rodriguez, P. (2024).
\newblock Benchmarking large language models on controllable generation under diversified constraints.
\newblock In \emph{AAAI}.

\bibitem[Liang et al.(2023)]{liang2023holistic}
Liang, P., Bommasani, R., Lee, T., et~al. (2023).
\newblock Holistic evaluation of language models.
\newblock \emph{Transactions on Machine Learning Research}.

\bibitem[Magar and Schwartz(2022)]{magar2022data}
Magar, I. and Schwartz, R. (2022).
\newblock Data contamination: From memorization to exploitation.
\newblock In \emph{ACL}.

\bibitem[Maris and Bechger(2009)]{maris2009equivalence}
Maris, G. and Bechger, T. (2009).
\newblock On interpreting the model parameters for the three parameter logistic model.
\newblock \emph{Measurement}, 7(2):75--88.

\bibitem[Mislevy(1986)]{mislevy1986bayes}
Mislevy, R.~J. (1986).
\newblock Bayes modal estimation in item response models.
\newblock \emph{Psychometrika}, 51(2):177--195.

\bibitem[Oren et al.(2024)]{oren2024proving}
Oren, Y., Meister, N., Chatterji, N., Ladhak, F., and Hashimoto, T. (2024).
\newblock Proving test set contamination in black box language models.
\newblock In \emph{EMNLP}.

\bibitem[Park et al.(2023)]{park2023trak}
Park, S.~M., Georgiev, K., Ilyas, A., Leclerc, G., and Madry, A. (2023).
\newblock {TRAK}: Attributing model behavior at scale.
\newblock In \emph{ICML}.

\bibitem[Platt(1999)]{platt1999probabilistic}
Platt, J. (1999).
\newblock Probabilistic outputs for support vector machines and comparisons to regularized likelihood methods.
\newblock \emph{Advances in Large Margin Classifiers}, 10(3):61--74.

\bibitem[Polo et al.(2024)]{polo2024tinybenchmarks}
Polo, F.~M., Choshen, L., Sun, W., et~al. (2024).
\newblock {tinyBenchmarks}: Evaluating {LLM}s with fewer examples.
\newblock In \emph{ICML}.

\bibitem[Pruthi et al.(2020)]{pruthi2020tracin}
Pruthi, G., Liu, F., Kale, S., and Sundararajan, M. (2020).
\newblock Estimating training data influence by tracing gradient descent.
\newblock In \emph{NeurIPS}.

\bibitem[Reckase(2009)]{reckase2009multidimensional}
Reckase, M.~D. (2009).
\newblock \emph{Multidimensional Item Response Theory}.
\newblock Springer.

\bibitem[Rodriguez et al.(2021)]{rodriguez2021evaluation}
Rodriguez, P., Barrow, J., Hoover, A.~M., Gribkoff, J., and Boyd-Graber, J. (2021).
\newblock Evaluation examples are not equally informative: How should that change {NLP} leaderboards?
\newblock In \emph{ACL}.

\bibitem[San~Mart\'in et al.(2009)]{sanmartin2009identifiability}
San~Mart\'in, E., Gonz\'alez, J., and Tuerlinckx, F. (2009).
\newblock On the identifiability of the fixed effects 3PL model.
\newblock \emph{Psychometrika}, 74(3):479--500.

\bibitem[Shi et al.(2024)]{shi2024detecting}
Shi, W., Ajith, A., Xia, M., et~al. (2024).
\newblock Detecting pretraining data from large language models.
\newblock In \emph{ICLR}.

\bibitem[Shokri et al.(2017)]{shokri2017membership}
Shokri, R., Stronati, M., Song, C., and Shmatikov, V. (2017).
\newblock Membership inference attacks against machine learning models.
\newblock In \emph{IEEE S\&P}.

\bibitem[Sinharay(2017)]{sinharay2017assessment}
Sinharay, S. (2017).
\newblock An assessment of the item compromise effect.
\newblock \emph{Journal of Educational Measurement}, 54(4):389--411.

\bibitem[Talts et al.(2018)]{talts2018validating}
Talts, S., Betancourt, M., Simpson, D., Vehtari, A., and Gelman, A. (2018).
\newblock Validating {B}ayesian inference algorithms with simulation-based calibration.
\newblock \emph{arXiv preprint arXiv:1804.06788}.

\bibitem[Thissen et al.(1993)]{thissen1993detection}
Thissen, D., Steinberg, L., and Wainer, H. (1993).
\newblock Detection of differential item functioning using the parameters of item response models.
\newblock In Holland and Wainer, editors, \emph{Differential Item Functioning}, pages 67--113.

\bibitem[White et al.(2024)]{white2024livebench}
White, C., Dooley, S., Roberts, M., Pal, A., and Goldblum, M. (2024).
\newblock {LiveBench}: A challenging, contamination-free {LLM} benchmark.
\newblock \emph{arXiv preprint arXiv:2406.19314}.

\bibitem[Xu et al.(2024)]{xu2024benchmarking}
Xu, L., Qi, A., Zhao, Y., and Chen, L. (2024).
\newblock Benchmarking benchmark leakage in large language models.
\newblock In \emph{ICML}.

\bibitem[Yang et al.(2023)]{yang2023rethinking}
Yang, S., Chiang, W.-L., Zheng, L., et~al. (2023).
\newblock Rethinking benchmark and contamination for language models with rephrased samples.
\newblock \emph{arXiv preprint arXiv:2311.04850}.

\bibitem[Zheng et al.(2023)]{zheng2023judging}
Zheng, L., Chiang, W.-L., Sheng, Y., et~al. (2023).
\newblock Judging {LLM}-as-a-judge with {MT-Bench} and {Chatbot Arena}.
\newblock In \emph{NeurIPS}.

\bibitem[Zhuang et al.(2024)]{zhuang2024hibayes}
Zhuang, Y., Zhang, Y., and Chen, J. (2024).
\newblock {HiBayES}: Hierarchical {B}ayesian evaluation of language models.
\newblock \emph{arXiv preprint}.

\bibitem[Zumbo(1999)]{zumbo1999handbook}
Zumbo, B.~D. (1999).
\newblock \emph{A Handbook on the Theory and Methods of Differential Item Functioning ({DIF})}.
\newblock National Defense Headquarters.

\end{thebibliography}

\end{document}
